\section{Statistical Evaluation Methodology}
\label{sec:statistical_methodology}

Evaluating high-stakes clinical decision support architectures requires rigorous statistical metrics beyond raw point-estimate accuracy. MedRAGAgents incorporates an automated statistical evaluation module (\texttt{evaluate.py}) that computes accuracy gains, output invalidity rates, non-parametric bootstrap confidence intervals, and paired significance tests.

\subsection{Primary Evaluation Metrics}

Let $N$ denote the total number of test vignettes evaluated in a given dataset. For each vignette $i \in \{1, \dots, N\}$, let $y_i^*$ represent the gold-standard option letter and $\hat{y}_i$ represent the system's predicted answer.

\begin{definition}[Classification Accuracy]
The classification accuracy $\text{Acc}(V_k)$ for system variant $V_k$ is the proportion of correctly predicted vignettes:
\begin{equation}
    \text{Acc}(V_k) = \frac{1}{N} \sum_{i=1}^{N} \mathbb{I}(\hat{y}_i = y_i^*)
\end{equation}
where $\mathbb{I}(\cdot)$ is the indicator function.
\end{definition}

\begin{definition}[Invalid Response Rate]
The invalid response rate $\text{IRR}(V_k)$ measures the frequency of outputs failing to yield a valid option letter $A \in \{A, B, C, D, E\}$:
\begin{equation}
    \text{IRR}(V_k) = \frac{1}{N} \sum_{i=1}^{N} \mathbb{I}(\hat{y}_i \notin \{A, B, C, D, E\})
\end{equation}
\end{definition}

\begin{definition}[Accuracy Gain]
The net accuracy gain $\Delta \text{Acc}(V_3, V_0)$ evaluates the marginal improvement of the full system (V3) over the unassisted baseline baseline (V0):
\begin{equation}
    \Delta \text{Acc}(V_3, V_0) = \text{Acc}(V_3) - \text{Acc}(V_0)
\end{equation}
\end{definition}

\subsection{Non-Parametric Bootstrap Confidence Intervals}

To account for variance over small test samples, 95\% confidence intervals are estimated using non-parametric bootstrap resampling with $B = 1,000$ iterations.

Let $\mathbf{D} = \{(y_1^*, \hat{y}_1), \dots, (y_N^*, \hat{y}_N)\}$ denote the observed dataset predictions. At each bootstrap iteration $b \in \{1, \dots, B\}$:
\begin{enumerate}
    \item Sample $N$ instances uniformly with replacement from $\mathbf{D}$, yielding bootstrap dataset $\mathbf{D}^{(b)}$.
    \item Compute the sample accuracy $\theta^{(b)} = \text{Acc}(\mathbf{D}^{(b)})$.
\end{enumerate}

The 95\% confidence interval $[L_{0.025}, U_{0.975}]$ is derived from the empirical percentiles of the ordered bootstrap distribution $\{\theta^{(1)}, \dots, \theta^{(B)}\}$:

\begin{equation}
    L_{0.025} = \text{Percentile}\left( \{\theta^{(b)}\}, 2.5\% \right), \quad U_{0.975} = \text{Percentile}\left( \{\theta^{(b)}\}, 97.5\% \right)
\end{equation}

\subsection{McNemar Significance Testing}

To verify whether performance improvements between baseline $V_0$ and full system $V_3$ are statistically significant, we construct a paired $2 \times 2$ outcome contingency matrix:

\begin{table}[h!]
\centering
\caption{Paired Outcome Contingency Matrix for McNemar Test.}
\label{tab:mcnemar_table}
\small
\begin{tabular}{c c c c}
\toprule
 & & \multicolumn{2}{c}{\textbf{Variant V3 Outcome}} \\
 & & \textbf{Correct ($\hat{y}_3 = y^*$)} & \textbf{Incorrect ($\hat{y}_3 \neq y^*$)} \\
\midrule
\multirow{2}{*}{\textbf{Variant V0 Outcome}} & \textbf{Correct ($\hat{y}_0 = y^*$)} & $a$ (Tie Correct) & $b$ (Loss for V3) \\
 & \textbf{Incorrect ($\hat{y}_0 \neq y^*$)} & $c$ (Win for V3) & $d$ (Tie Incorrect) \\
\bottomrule
\end{tabular}
\end{table}

The McNemar test statistic $\chi^2$ assesses the null hypothesis $H_0: P(b) = P(c)$ of equal discordant error rates:

\begin{equation}
    \chi^2 = \frac{(|c - b| - 1)^2}{c + b}
\end{equation}

Under $H_0$, the statistic $\chi^2$ follows a chi-squared distribution with 1 degree of freedom, yielding the upper-tail $p$-value:
\begin{equation}
    p\text{-value} = P\left( \chi_1^2 \ge \chi^2 \right)
\end{equation}
